\section{Related Work}

\subsection{Reliability signals for fully test-time adaptation}

Fully test-time adaptation often starts from entropy minimization. \textsc{Tent} updates normalization parameters by minimizing prediction entropy on unlabeled target data~\citep{wang2021tent}, and MEMO-style objectives add augmentation consistency~\citep{bartler2022memo}; related variants use contrastive objectives or parameter-free online updates~\citep{chen2022ctta,boudiaf2022potta}. Later work shows that low entropy alone is incomplete: \textsc{EATA} removes high-entropy and redundant samples while preserving source-important parameters~\citep{niu2022eata}, \textsc{SAR} combines reliable entropy with sharpness-aware recovery~\citep{niu2023sar}, \textsc{DeYO} uses destructive transformations to expose fragile confidence~\citep{lee2024deyo}, and \textsc{AdaDEM} analyzes reward collapse and easy-class bias~\citep{ma2025adadem}. \method{} follows this diagnosis, but \core{} shifts the reliability unit from a scalar sample score to a corroborated tensor entity checked by source, raw, photometric, structural, and destroy-view evidence.

\subsection{Drift control in continual and wild streams}

Continual TTA exposes cumulative drift because each accepted pseudo-supervised update changes the model that will supervise later inputs. \textsc{CoTTA} uses a mean teacher, augmentation-averaged pseudo-labels, and stochastic restoration~\citep{wang2022cotta}; \textsc{RoTTA} adds robust batch normalization and category-balanced memory~\citep{yuan2023rotta}; NOTE, DSS, and related teacher or memory methods stabilize online adaptation through replay, selection, or self-distillation~\citep{gong2022note,hong2024dss,song2023ecotta,doebler2023rmt}. PALM studies adaptive learning-rate control in continual TTA~\citep{maharana2025palm}. Lifelong Test-Time Adaptation constrains adaptation by projecting gradients into a tracked low-dimensional subspace~\citep{duan2025lifelong}, connecting continual TTA to intrinsic-dimension, low-dimensional training, and historical-solution averaging analyses~\citep{larsen2021intrinsic,li2022lowdim,li2023twa}. These results show that history and anchors matter in non-i.i.d. streams, but teachers, memories, normalization statistics, and gradient bases remain tied to specific stream and architecture assumptions. \care{} instead puts the anchor inside the loss for each accepted entity, using source-relative clipping or class-centered support rather than a separate stateful safeguard.

\subsection{Prompt and dense-prediction adaptation}

Vision-language and dense-prediction TTA broaden the adaptation entity beyond image-level classification. Prompt-only adaptation also draws on learned prompts, where CoOp, CoCoOp, ProDA, and MaPLe establish strong parameterizations for CLIP-style models~\citep{zhou2022learning,zhou2022conditional,khattak2022proda,khattak2023maple}. TPT adapts prompts through entropy minimization over selected augmented views~\citep{shu2022tpt}, while SwapPrompt, diffusion-enhanced TTA, DPCore, ZERO, TDA, FreeTTA, OT-VP, and FOA study prompt replacement, text-and-image augmentation, continual prompt memory, multi-view marginalization, caches, online distribution modeling, optimal transport, and forward-only VLM adaptation~\citep{ma2023swapprompt,feng2025diffusion,zhang2024dpcore,farina2024zero,karmanov2024tda,dai2025freetta,zhang2025otvp,niu2024foa}. Dense prediction introduces pixel-level imbalance, so segmentation adaptation studies domain properties, open-vocabulary optimization, correlation alignment, or hybrid adaptation~\citep{jeon2025exploiting,de2025test,you2025test,park2025hybrid}. These designs expose the same interface problem: images, pixels, and prompt views produce losses at different scales. \sane{} addresses it by normalizing \atlasloss{} over active selected entities, so update magnitude depends on accepted evidence rather than raw tensor count.
